\documentclass[11pt]{article}
\usepackage[margin=.85in]{geometry}
\usepackage{amsmath,amssymb}
\usepackage[hidelinks]{hyperref}

\title{Conjecture 00000000457: Shi coefficient log-concavity}
\author{}
\date{}

\begin{document}
\maketitle

For the Shi arrangement in $\mathbb R^n$ with hyperplanes
$x_i-x_j=0,1$ for $1\le i<j\le n$, Headley's formula gives
$\chi(t)=t(t-n)^{n-1}$ \cite[Section~1]{athanasiadis}.
Restricting to the sum-zero hyperplane removes the common diagonal
line. This leaves the same intersection poset and lowers the
dimension of every intersection by one, so the essential arrangement
has characteristic polynomial $(t-n)^{n-1}$.
Write $N=n-1$. In descending degree order, the absolute values of the
nonzero coefficients in either convention are
\[
 a_k=\binom{N}{k}n^k,\qquad 0\le k\le N.
\]

Use the circular parking model with $n$ spots and available car
labels $1,\ldots,N$. An element of $\mathcal P_k$ is a set of $k$
labels with one preferred spot assigned to each. Cars arrive in
increasing label order and advance cyclically from their preferred
spot until finding an empty one. Each car parks because fewer than
$n$ spots have already been occupied. Hence every choice of labels
and preferences is a successful circular parking word, and
$|\mathcal P_k|=a_k$.
Any spot empty at the end was empty throughout the process. Cutting
the circle immediately after such a spot gives a linear
parking word on the remaining spots, since no car passed the chosen
empty spot.

Fix $1\le k<N$ and take a pair in
$\mathcal P_{k-1}\times\mathcal P_{k+1}$ with label sets $A,B$.
Read the labels of $A\mathbin{\triangle}B$ in increasing order.
The running excess of $B$-only labels over $A$-only labels starts
at zero, changes by $\pm1$ at each step, and ends at $2$. Stop at
the first prefix where it equals $1$. Exchange the membership of
every label in this prefix, moving each car with its preferred
spot. Keep the two preferences of each common label in their
original words. Exactly one more car moves from $B$ to $A$ than
from $A$ to $B$, yielding a pair in
$\mathcal P_k\times\mathcal P_k$.

The exchange preserves the ordered symmetric difference. Within
the exchanged prefix, the new excess of first-word labels over
second-word labels equals the old excess of $B$-only labels over
$A$-only labels. The same stopping prefix is therefore recovered
as the first prefix reaching $1$ for this new excess. Exchanging
it again recovers both input words, so the map is injective.
A target pair with identical label sets has empty symmetric
difference and cannot be in the image. Thus $a_k^2>a_{k-1}a_{k+1}$
throughout the interior of the positive support.

The positive ratios $a_k/a_{k-1}$ consequently decrease, so the
sequence rises until those ratios fall below $1$ and falls after
that, allowing a tie at its maximum. The full polynomial appends
a zero constant coefficient, which preserves both log-concavity
and unimodality. The cases with no interior positive index are
immediate from the displayed sequence. This proves the coefficient
inequality and the resulting unimodality by a parking-word injection.

The Lean theorem \texttt{exchange\_certificate} checks that the
prefix exchange balances the exclusive labels, carries their payloads,
has the stated inverse, and leaves a nonempty symmetric difference.
The theorem \texttt{exchange\_parking\_sizes} restores the common
labels to give exactly $k$ cars in each output word. The characteristic
polynomial formula and the parking-word enumeration are used above.

\begin{thebibliography}{9}
\bibitem{athanasiadis}
C. A. Athanasiadis, \emph{Inductively Free Arrangements and a Result
of Headley}, FPSAC, 1997.
\href{https://fpsac-archive.github.io/FPSAC97/ARTICLES/Athanasiadis.pdf}{Proceedings paper}.
\end{thebibliography}

\end{document}
